\documentclass[../../../include/open-logic-section]{subfiles}

\begin{document}

\olfileid[fr]{his}{set}{hilbertcurve}

\olsection{Annexe : les courbes de Hilbert remplissant l'espace}

Dans la \olref[pathology]{sec}, nous avons indiqué que
la démonstration de Cantor selon laquelle une ligne et un carré
ont exactement le même nombre de points
(\olref[his][set][cantorplane]{thm:cantorplane}) incita Peano à
chercher une \emph{courbe} remplissant l'espace. Il en construisit
une en 1890. Nous allons expliquer, de manière informelle, comment
construire la courbe de Hilbert, avec une petite modification.\footnote{
Pour une explication plus rigoureuse, voir \cite{Rose2010}. La
modification consiste à ajouter les parties rouges des courbes
ci-dessous, ce qui facilite un peu la vérification de leur
continuité.}\footnote{Corrections mathématiques de l'édition
française : la source donne à la variable de la limite le carré
pour domaine, alors qu'il faut le segment. Elle ne précise pas
la paramétrisation commune des approximations et son argument
final ne démontre pas la continuité en chaque point du domaine.
Nous rendons ici explicites la paramétrisation, la convergence
uniforme et l'argument de surjectivité.}

Nous devons définir une application $h$ comme limite d'une suite
d'applications $h_1$, $h_2$, $h_3$, \dots\@ Nous décrirons d'abord
la construction, puis montrerons que son image remplit le carré
et qu'elle définit bien une courbe.

L'image recherchée est le carré unité $\unitsquare$. Voici la
première approximation de $h$, c'est-à-dire $h_1$ :
\begin{center}
	\begin{tikzpicture}
	\draw[gray] (0pt, 0pt) rectangle (80pt, 80pt);
	\draw[step = 40pt, gray, very thin] (0pt, 0pt) grid (80pt, 80pt);
	\draw[oldiagcolorC, thick] (20pt, 0)--(20pt, 20pt);
	\draw[oldiagcolorC, thick] (60pt, 0)--(60pt, 20pt);
	\begin{scope}[xshift=20pt, yshift=20pt]
	\draw [black, thick, l-system={Hilbert curve, step=40pt, angle=90, axiom=L, order=1}]  lindenmayer system;
	\end{scope}
	\end{tikzpicture}
\end{center}
Pour suivre la construction, nous avons superposé au carré une
grille $2 \times 2$. La courbe commence dans le quart inférieur
gauche, passe au quart supérieur gauche, puis au quart supérieur
droit, et se termine dans le quart inférieur droit. Voici la
deuxième étape, c'est-à-dire $h_2$ :
\begin{center}
	\begin{tikzpicture}
	\draw[gray] (0pt, 0pt) rectangle (80pt, 80pt);
	\draw[step = 20pt, gray, very thin] (0pt, 0pt) grid (80pt, 80pt);
	\draw[oldiagcolorC, thick] (0pt, 10pt)--(10pt, 10pt);
	\draw[oldiagcolorC, thick] (70pt, 10pt)--(80pt, 10pt);
	\draw[thick, oldiagcolorE] (10pt, 30pt)--(10pt, 50pt);
	\draw[thick, oldiagcolorE] (30pt, 50pt)--(50pt, 50pt);
	\draw[thick, oldiagcolorE] (70pt, 50pt)--(70pt, 30pt); \begin{scope}[xshift=10pt, yshift=30pt]
	\draw [thick, rotate=270,black, l-system={Hilbert curve, step=20pt, angle=90, axiom=L, order=1}]  lindenmayer system;
	\end{scope}
	\begin{scope}[xshift=10pt, yshift=50pt]
	\draw [thick, black, l-system={Hilbert curve, step=20pt, angle=90, axiom=L, order=1}]  lindenmayer system;
	\end{scope}
	\begin{scope}[xshift=50pt, yshift=50pt]
	\draw [thick, black, l-system={Hilbert curve, step=20pt, angle=90, axiom=L, order=1}]  lindenmayer system;
	\end{scope}
	\begin{scope}[xshift=70pt, yshift=10pt]
	\draw [thick, rotate=90,black, l-system={Hilbert curve, step=20pt, angle=90, axiom=L, order=1}]  lindenmayer system;
	\end{scope}
	\end{tikzpicture}
\end{center}
Les couleurs expliquent la construction de $h_2$. Nous plaçons
d'abord des copies réduites de la partie qui n'est pas
!!{colorC} dans $h_1$ aux quatre quarts du carré, dans le même
ordre : inférieur gauche, supérieur gauche, supérieur droit,
inférieur droit. Ces copies sont dessinées en noir. Nous les
relions ensuite par des segments !!{colorE}s, puis relions la
figure au bord du carré par des segments !!{colorC}s.

Passons à $h_3$. De même que $h_2$ était formée de quatre copies
réduites et reliées de la partie non rouge de $h_1$, $h_3$ est
formée de quatre copies réduites de la partie non rouge de $h_2$,
dessinées en noir. Elles sont reliées par des segments !!{colorE}s,
puis raccordées au bord du carré par des segments !!{colorC}s.
\begin{center}
	\begin{tikzpicture}
	\draw[gray] (0pt, 0pt) rectangle (80pt, 80pt);
	\draw[step = 10pt, gray, very thin] (0pt, 0pt) grid (80pt, 80pt);
	\draw[thick, oldiagcolorC](5pt, 0pt)--(5pt, 5pt);
	\draw[thick, oldiagcolorC] (75pt, 0pt)--(75pt, 5pt);
	\draw[thick, oldiagcolorE] (75pt, 45pt)--(75pt, 35pt);
	\draw[thick, oldiagcolorE] (5pt, 35pt)--(5pt, 45pt);
	\draw[thick, oldiagcolorE] (35pt, 45pt)--(45pt, 45pt);
	\draw[thick, oldiagcolorE] (75pt, 45pt)--(75pt, 35pt); \begin{scope}[xshift=5pt, yshift=35pt]
	\draw [rotate=270, thick, black, l-system={Hilbert curve, step=10pt, angle=90, axiom=L, order=2}]  lindenmayer system;
	\end{scope}
	\begin{scope}[xshift=5pt, yshift=45pt]
	\draw [thick, black, l-system={Hilbert curve, step=10pt, angle=90, axiom=L, order=2}]  lindenmayer system;
	\end{scope}
	\begin{scope}[xshift=45pt, yshift=45pt]
	\draw [thick, black, l-system={Hilbert curve, step=10pt, angle=90, axiom=L, order=2}]  lindenmayer system;
	\end{scope}
	\begin{scope}[xshift=75pt, yshift=5pt]
	\draw [thick, rotate=90,black, l-system={Hilbert curve, step=10pt, angle=90, axiom=L, order=2}]  lindenmayer system;
	\end{scope}
	\end{tikzpicture}
\end{center}
Le procédé général permettant de construire $h_{n+1}$ à partir
de $h_n$ apparaît maintenant.

Précisons comment ces dessins deviennent des applications définies
sur le même segment $\unitline=[0,1]$. À l'étape $n$, le parcours
visite successivement les $4^n$ cellules fermées de la grille
$2^n \times 2^n$. Partageons $[0,1]$ en $4^n$ sous-intervalles
égaux, dans cet ordre, et faisons parcourir de façon continue à
$h_n$ la partie du tracé située dans la cellule correspondante
pendant chacun de ces sous-intervalles. Les raccords entre
cellules adjacentes assurent la continuité de $h_n$.
Le raffinement partage ensuite chaque sous-intervalle en quatre
et chacune des cellules correspondantes en quatre, selon les
orientations indiquées par la construction. Ainsi, si $I$ est
un sous-intervalle de l'étape $n$ associé à une cellule $Q$,
toutes les étapes $m\geq n$ vérifient $h_m(I)\subseteq Q$.
Aux extrémités communes des sous-intervalles, les raccords se
trouvent sur les bords communs des cellules.

Par le théorème de Pythagore, le diamètre de chaque cellule vaut :
\[
\sqrt{\left(\nicefrac{1}{2^{n}}\right)^2+\left(\nicefrac{1}{2^{n}}\right)^2} = 2^{(\frac{1}{2}-n)}.
\]
Pour un même $x\in\unitline$, les points $h_m(x)$ et $h_k(x)$
appartiennent donc à une même cellule de l'étape $n$ dès que
$m,k\geq n$. Leur distance est au plus
$2^{(\frac{1}{2}-n)}$, indépendamment de~$x$. Les suites
$h_n(x)$ sont ainsi de Cauchy, uniformément en~$x$, et convergent
dans le carré fermé. Nous pouvons enfin définir la courbe
\emph{elle-même}, pour chaque $x \in \unitline$, par :
\begin{align*}
	h(x) &= \lim_{n \rightarrow \infty} h_n(x).
\end{align*}
La même estimation donne
$\|h(x)-h_n(x)\|\leq 2^{(\frac{1}{2}-n)}$ pour tout~$x$.
La convergence est donc uniforme.

Montrons que l'image de $h$ est tout le carré. Chaque tracé
$h_n$ traverse toutes les cellules de la grille ; tout point du
carré se trouve donc à une distance \emph{au plus} égale à
$2^{(\frac{1}{2}-n)}$ d'un point de ce tracé. Fixons une cellule
$Q$ de l'étape~$n$ et le sous-intervalle~$I$ qui lui correspond.
Comme $Q$ est fermé, $h(I)\subseteq Q$. Montrons l'inclusion inverse.
Pour tout point $p\in Q$, choisissons
une suite de cellules fermées emboîtées $Q_k$ contenant~$p$,
en commençant par $Q_n=Q$ et en prenant à chaque étape une
des quatre sous-cellules. Les sous-intervalles fermés
correspondants $I_k$ sont emboîtés, avec $I_n=I$, et ont pour
longueurs $4^{-k}$. Leur intersection contient un unique
point~$x$. Pour tout $k\geq n$, toutes les approximations
d'ordre $m\geq k$ envoient $x$ dans $Q_k$, qui contient
donc aussi leur limite $h(x)$. Comme le diamètre de $Q_k$
tend vers zéro, $h(x)=p$. Ainsi, $h(I)=Q$ et chaque point du carré
se trouve bien \emph{sur} la courbe. La distance maximale des
points du carré à l'image de $h$ est donc nulle, conformément à :
\[
\lim_{n \rightarrow \infty} 2^{(\frac{1}{2}-n)} = 0.
\]
La courbe $h$ remplit le carré !

Il reste à vérifier que $h$ est bien une \emph{courbe}, au sens
précis de ce mot. La définition moderne s'appuie sur celle donnée
par Jordan en 1887, quelques années seulement avant la première
courbe remplissant l'espace :

\begin{defn}
Une courbe est une application continue de $\unitline$ dans
$\Real^2$.
\end{defn}

Intuitivement, une courbe décrit un parcours continu dans le plan,
paramétré par un segment. La continuité ne signifie pas que le
parcours est lisse ou dérivable. Notre application $h$ a déjà
pour domaine $\unitline$ et pour codomaine $\Real^2$ ; il reste
à établir sa continuité.

Nous avons défini la continuité dans la \olref[limits]{sec} à l'aide
de la notation $\epsilon$/$\delta$. Fixons donc un point
$x_0\in\unitline$ et une précision $\epsilon>0$, et posons
$p=h(x_0)$. Il faut trouver $\delta>0$ tel que tous les
$x\in\unitline$ satisfaisant $|x-x_0|<\delta$ aient leur image
dans le disque de centre $p$ et de rayon $\epsilon$. Le disque
peut déborder du carré sans que cela pose de problème.

Choisissons $n$ assez grand pour que
$2^{(\frac{1}{2}-n)}<\epsilon/3$. Puisque $h_n$ est continue,
il existe $\delta>0$ tel que, pour tout $x\in\unitline$,
$|x-x_0|<\delta$ implique
$\|h_n(x)-h_n(x_0)\|<\epsilon/3$. L'estimation uniforme donne
alors :
\[
\|h(x)-h(x_0)\|
\leq \|h(x)-h_n(x)\|
+\|h_n(x)-h_n(x_0)\|
+\|h_n(x_0)-h(x_0)\|
<\epsilon.
\]
C'est exactement la continuité de $h$ en~$x_0$. Le point
$x_0$ étant arbitraire, $h$ est continue sur $\unitline$.
Nous avons donc construit une courbe qui remplit le carré unité.

\end{document}
